\chapter{Pediatric and Neonatal Diseases}
\label{ch:pediatric}
Pediatric and neonatal diseases encompass a broad spectrum of conditions affecting infants, children, and adolescents. This chapter covers prematurity and its complications, perinatal conditions, congenital anomalies, genetic and chromosomal disorders, common childhood illnesses, and pediatric cancers. Many of these conditions are unique to the developing child or have distinct presentations, diagnostic considerations, and management approaches compared to adult diseases.

%-------------------------------------------------------------
\section{Common Childhood Illnesses}
%-------------------------------------------------------------

%-------------------------------------------------------------

%-------------------------------------------------------------

%-------------------------------------------------------------

%-------------------------------------------------------------

%-------------------------------------------------------------

%-------------------------------------------------------------

%-------------------------------------------------------------

\label{sec:Common Childhood Illnesses}
%-------------------------------------------------------------

\subsection{Neural Tube Defects}
\index{Neural Tube Defects}
\label{sec:Neural Tube Defects}
Neural tube defects (NTDs) are congenital malformations resulting from failure of neural tube closure during the fourth week of embryonic development, classified as open or closed. Incidence is 0.5--1 per 1000 pregnancies; risk is reduced by 50--70\% with periconceptional folic acid supplementation. Spina bifida encompasses a spectrum from spina bifida occulta (vertebral involvement only), meningocele (meningeal herniation), to meningomyelocele (herniation of meninges and spinal cord tissue causing neurologic deficits). Anencephaly is a fatal NTD with absence of the cerebral hemispheres. Diagnosis is by prenatal ultrasound and maternal serum alpha-fetoprotein screening. Management of meningomyelocele includes prenatal or postnatal surgical repair and multidisciplinary care. Prognosis depends on the level and severity of the lesion.

\subsection{Neuroblastoma}
\index{Neuroblastoma}
\label{sec:Neuroblastoma}
Neuroblastoma is the most common extracranial solid tumor in children and the most common malignancy in infants. It arises from the sympathetic nervous system, most commonly in the adrenal medulla (40\%) or paraspinal sympathetic chain. Incidence is 1 in 7000 children; median age at diagnosis is 18 months. Clinical features depend on the site: abdominal mass (adrenal), Horner syndrome (cervical), opsoclonus-myoclonus-ataxia (paraneoplastic), proptosis and periorbital ecchymoses (orbital metastases), and bone pain (skeletal metastases). Diagnosis is by biopsy with histology, MYCN amplification status, and imaging. International Neuroblastoma Risk Group staging classifies tumors as low-risk, intermediate-risk, or high-risk. Treatment varies by risk group: low-risk tumors may be observed or surgically resected; intermediate-risk tumors receive chemotherapy; high-risk tumors require intensive multimodal therapy. Prognosis ranges from excellent (>95\% survival in low-risk) to poor (40--50\% survival in high-risk).

\subsection{Osteosarcoma in Children}
\index{Osteosarcoma in Children}
\label{sec:Osteosarcoma in Children}
Osteosarcoma is the most common primary malignant bone tumor in children and adolescents (see \hyperref[ch:bones]{Chapter~26, Diseases of the Musculoskeletal System} for comprehensive coverage). Peak incidence is during the adolescent growth spurt (ages 10--16), with a predilection for the metaphyses of long bones, most commonly the distal femur, proximal tibia, and proximal humerus. Most childhood osteosarcomas are high-grade. Treatment follows neoadjuvant chemotherapy, surgical resection, and adjuvant chemotherapy, with survival of 60--70\% for localized disease.

\subsection{Patau Syndrome (Trisomy 13)}
\index{Patau Syndrome (Trisomy 13)}
\label{sec:Patau Syndrome}
Patau syndrome is a chromosomal disorder caused by trisomy 13, with an incidence of 1 in 10,000 live births. Clinical features include severe intellectual disability, microcephaly, holoprosencephaly, cleft lip and palate, polydactyly, low-set ears, microphthalmia, and congenital heart disease. Median survival is 3--10 days; less than 10\% survive to age 1 year. Diagnosis is by karyotype. Management is supportive and palliative.

\subsection{Patent Ductus Arteriosus (PDA) in Prematurity}
\index{Patent Ductus Arteriosus (PDA) in Prematurity}
\label{sec:Patent Ductus Arteriosus in Prematurity}
Patent ductus arteriosus (PDA) is the persistent patency of the fetal vascular connection between the descending aorta and the pulmonary artery. In preterm infants, ductal closure is delayed, and a hemodynamically significant PDA occurs in up to 60\% of extremely preterm infants. Left-to-right shunting through the PDA leads to pulmonary overcirculation, increased left ventricular preload, and systemic steal. Clinical features include a continuous murmur, bounding pulses, wide pulse pressure, respiratory deterioration, and hypotension. Diagnosis is by echocardiography. Treatment options include fluid restriction, diuretics, pharmacologic closure with indomethacin or ibuprofen, and surgical ligation for failed medical therapy. Prognosis is favorable with closure.

\subsection{Pediatric Appendicitis}
\index{Pediatric Appendicitis}
\label{sec:Pediatric Appendicitis}
Appendicitis is the most common surgical emergency in children (see \hyperref[ch:intestines]{Chapter~16, Diseases of the Intestines} for comprehensive coverage). In pediatric patients, diagnosis may be more challenging due to atypical presentations, especially in younger children. Clinical features include periumbilical pain migrating to the right lower quadrant, fever, anorexia, nausea and vomiting, and rebound tenderness. Ultrasound is the first-line imaging modality. Treatment is laparoscopic appendectomy with perioperative antibiotics. Prognosis is excellent with timely treatment; young children have higher rates of perforation.

\subsection{Perinatal Asphyxia / Hypoxic-Ischemic Encephalopathy (HIE)}
\index{Perinatal Asphyxia / Hypoxic-Ischemic Encephalopathy (HIE)}
\label{sec:Hypoxic-Ischemic Encephalopathy}
Hypoxic-ischemic encephalopathy (HIE) is brain injury resulting from perinatal asphyxia---impaired cerebral blood flow and oxygen delivery around the time of birth. Incidence is 1--3 per 1000 live births in developed countries. The condition results from primary energy failure (decreased ATP, cellular necrosis) followed by a secondary phase of inflammation, excitotoxicity, and apoptosis hours to days later. Clinical features include low Apgar scores, metabolic acidosis, seizures, hypotonia, abnormal neurological examination, and multi-organ dysfunction. HIE is graded as mild (Sarnat Stage 1), moderate (Stage 2), or severe (Stage 3). Diagnosis is based on perinatal history, blood gas analysis, and amplitude-integrated EEG; MRI within the first week shows characteristic patterns of injury. Treatment is therapeutic hypothermia (33.5$^\circ$C for 72 hours), initiated within 6 hours of birth, which reduces mortality and neurodevelopmental disability. Prognosis depends on severity and timeliness of treatment.

\subsection{Periventricular Leukomalacia (PVL)}
\index{Periventricular Leukomalacia (PVL)}
\label{sec:Periventricular Leukomalacia}
Periventricular leukomalacia (PVL) is ischemic necrosis of the periventricular white matter in preterm infants, resulting from the vulnerability of oligodendrocyte precursor cells to hypoxic-ischemic injury. It occurs in 3--10\% of preterm infants and is a major cause of cerebral palsy, particularly spastic diplegia. The condition results from impaired perfusion in the watershed zone between the anterior, middle, and posterior cerebral arteries. Clinical features include neurodevelopmental impairment, with spastic diplegia as a major consequence. Diagnosis is by cranial ultrasound (cystic echolucencies) or MRI. Management focuses on prevention (avoiding hypotension, hypoxia, and hypocarbia). Prognosis is guarded; long-term outcomes include spastic diplegia, visual impairment, and cognitive deficits.

\subsection{Persistent Pulmonary Hypertension of the Newborn (PPHN)}
\index{Persistent Pulmonary Hypertension of the Newborn (PPHN)}
\label{sec:Persistent Pulmonary Hypertension of the Newborn}
Persistent pulmonary hypertension of the newborn (PPHN) is the failure of the pulmonary vascular resistance to decrease normally after birth, leading to right-to-left shunting across the foramen ovale and ductus arteriosus and profound hypoxemia. Incidence is approximately 2 per 1000 live births. The condition results from failure of the normal postnatal drop in pulmonary vascular resistance. Clinical features include severe hypoxemia out of proportion to the degree of lung disease, with a preductal-to-postductal oxygen saturation gradient. Diagnosis is by echocardiography demonstrating elevated pulmonary artery pressure and right-to-left shunt. Treatment includes optimization of lung inflation, sedation, inotropic support, inhaled nitric oxide, and ECMO for refractory cases. Prognosis depends on the underlying cause and severity.

\subsection{Prader-Willi Syndrome}
\index{Prader-Willi Syndrome}
\label{sec:Prader-Willi Syndrome}
Prader-Willi syndrome (PWS) is a neurogenetic disorder caused by loss of function of the paternal copy of genes on chromosome 15q11-q13, the region imprinted on the maternal allele. Mechanisms include paternal deletion (70\%), maternal uniparental disomy (25\%), and imprinting defects (5\%). Incidence is 1 in 10,000--30,000. Clinical features are biphasic: in infancy, hypotonia, poor feeding, and failure to thrive; in childhood, hyperphagia (insatiable appetite), obesity, short stature, hypogonadism, small hands and feet, and mild to moderate intellectual disability. Diagnosis is by methylation analysis. Management includes growth hormone therapy, strict dietary supervision to prevent obesity, and behavioral interventions.

\subsection{Preterm Birth}
\index{Preterm Birth}
\label{sec:Preterm Birth}
Preterm birth is a perinatal condition defined as delivery before 37 weeks of gestation. It affects approximately 10\% of all pregnancies worldwide and is the leading cause of neonatal morbidity and mortality. Risk factors include multiple gestation, intrauterine infection, cervical insufficiency, prior preterm birth, smoking, and maternal stress. The condition results from multiple pathogenic pathways including infection, inflammation, and uteroplacental ischemia. Clinical features include delivery before 37 weeks, classified as late preterm (34--36 weeks), moderate preterm (32--33 weeks), very preterm (28--31 weeks), and extremely preterm (<28 weeks), with complications inversely related to gestational age. Diagnosis is clinical. Management includes antenatal corticosteroids (betamethasone) to accelerate fetal lung maturation, magnesium sulfate for neuroprotection, and transfer to a facility with a level III NICU; tocolysis may be used to delay delivery for corticosteroid administration. Prognosis depends on gestational age and birth weight.

%-------------------------------------------------------------
\section{Congenital Anomalies}
%-------------------------------------------------------------

%-------------------------------------------------------------

%-------------------------------------------------------------

%-------------------------------------------------------------

%-------------------------------------------------------------

%-------------------------------------------------------------

%-------------------------------------------------------------

%-------------------------------------------------------------

\label{sec:Congenital Anomalies}
%-------------------------------------------------------------

\subsection{Down Syndrome (Trisomy 21)}
\index{Down Syndrome (Trisomy 21)}
\label{sec:Down Syndrome}
Down syndrome is the most common chromosomal disorder, with an incidence of 1 in 700 live births. It is caused by trisomy 21 (94\%), Robertsonian translocation (4\%), or mosaic trisomy 21 (2\%); risk increases with maternal age. Clinical features include flat facial profile, upslanting palpebral fissures, epicanthal folds, Brushfield spots, single transverse palmar crease, hypotonia, joint hyperlaxity, short stature, intellectual disability (mild to moderate), and congenital heart disease (50\%). Associated conditions include hearing loss, hypothyroidism, celiac disease, duodenal atresia, and increased risk of leukemia. Prenatal screening includes first-trimester nuchal translucency and maternal serum markers. Diagnosis is confirmed by karyotype or FISH. Management requires lifelong multidisciplinary care; life expectancy has increased to approximately 60 years.

\subsection{Edwards Syndrome (Trisomy 18)}
\index{Edwards Syndrome (Trisomy 18)}
\label{sec:Edwards Syndrome}
Edwards syndrome is a chromosomal disorder caused by trisomy 18, with an incidence of 1 in 5000 live births. Clinical features include severe intellectual disability, growth restriction, low-set malformed ears, micrognathia, prominent occiput, clenched fists with overlapping fingers, rocker-bottom feet, and congenital heart disease. Median survival is less than 1 year; only 5--10\% survive to age 1 year. Diagnosis is by karyotype. Management is supportive and palliative.

\subsection{Epiglottitis}
\index{Epiglottitis}
\label{sec:Epiglottitis}
Epiglottitis is a rapidly progressive, life-threatening inflammation of the epiglottis and supraglottic structures. Historically caused by Haemophilus influenzae type b (Hib), its incidence has declined dramatically due to the Hib vaccine. Clinical features include abrupt onset of high fever, severe sore throat, drooling, muffled voice (``hot potato'' voice), tripod positioning, and inspiratory stridor. Diagnosis is clinical; lateral neck radiograph shows the ``thumbprint sign'' (enlarged epiglottis). Airway management is the priority: immediate establishment of a secured airway (endotracheal intubation in the operating room). Blood and epiglottic cultures are obtained, and intravenous antibiotics (ceftriaxone or cefotaxime) are initiated. Prognosis is excellent with prompt airway management.

\subsection{Esophageal Atresia / Tracheoesophageal Fistula}
\index{Esophageal Atresia / Tracheoesophageal Fistula}
\label{sec:Esophageal Atresia}
Esophageal atresia (EA) with or without tracheoesophageal fistula (TEF) is a congenital anomaly with incidence of 1 per 3500 live births. The most common type (86\%) is EA with distal TEF (Gross Type C). Clinical presentation includes polyhydramnios (prenatal), excessive drooling, coughing, choking, and cyanosis with feeding. Failure to pass a nasogastric tube beyond 8--10 cm from the gums is diagnostic. Radiography shows a coiled nasogastric tube in the proximal esophageal pouch and gas in the stomach (if distal TEF is present). Associated anomalies include the VACTERL association. Treatment is surgical: primary repair with division of the fistula and end-to-end esophageal anastomosis. Prognosis is excellent for isolated EA/TEF.

\subsection{Ewing Sarcoma in Children}
\index{Ewing Sarcoma in Children}
\label{sec:Ewing Sarcoma in Children}
Ewing sarcoma is the second most common malignant bone tumor in children and adolescents (see \hyperref[ch:bones]{Chapter~26, Diseases of the Musculoskeletal System} for comprehensive coverage). Peak incidence is 10--20 years. It arises in the diaphysis of long bones and flat bones (pelvis, ribs). The hallmark is a t(11;22)(q24;q12) translocation producing the EWSR1-FLI1 fusion gene. Treatment includes chemotherapy, surgical resection, and/or radiation therapy. Survival for localized disease is 70--75\%.

\subsection{Fragile X Syndrome}
\index{Fragile X Syndrome}
\label{sec:Fragile X Syndrome}
Fragile X syndrome is the most common inherited cause of intellectual disability and the most common monogenic cause of autism spectrum disorder. It is caused by a CGG trinucleotide repeat expansion (>200 repeats) in the FMR1 gene on the X chromosome, leading to hypermethylation and gene silencing. Incidence is 1 in 4000 males and 1 in 8000 females. Clinical features in males include intellectual disability, long narrow face, prominent ears, large testes (macroorchidism, postpubertal), joint hypermobility, autistic behaviors, attention deficits, and anxiety. Females may have mild intellectual disability or learning difficulties. Premutation carriers (55--200 repeats) are at risk of FXTAS and FXPOI. Diagnosis is by FMR1 DNA testing. Treatment is supportive and includes behavioral therapies, educational interventions, and pharmacotherapy for associated conditions.

\subsection{Henoch-Sch\"onlein Purpura (IgA Vasculitis)}
\index{Henoch-Sch\"onlein Purpura (IgA Vasculitis)}
\label{sec:Henoch-Schonlein Purpura}
Henoch-Sch\"onlein purpura (HSP) is an IgA-mediated small-vessel vasculitis, the most common vasculitis in children (see \hyperref[ch:connective_tissue]{Chapter~27, Connective Tissue Diseases} for comprehensive coverage). Classic tetrad: palpable purpura (dependent areas), arthritis/arthralgia, abdominal pain, and renal involvement (hematuria, proteinuria). Mean age is 4--6 years. Diagnosis is clinical; skin or renal biopsy shows IgA deposition. Treatment is supportive; corticosteroids are used for severe abdominal pain or significant renal disease. Prognosis is excellent; most cases resolve spontaneously.

\subsection{Hepatoblastoma}
\index{Hepatoblastoma}
\label{sec:Hepatoblastoma}
Hepatoblastoma is the most common malignant liver tumor in children, primarily affecting children under 3 years of age. Clinical features include an asymptomatic abdominal mass, abdominal distension, and occasionally pain, anorexia, and jaundice. Serum alpha-fetoprotein is markedly elevated in >90\% of cases. Diagnosis is by abdominal ultrasound, CT/MRI, and AFP level. Complete surgical resection is the cornerstone of treatment, with neoadjuvant chemotherapy for initially unresectable tumors. Survival is approximately 70--80\% with complete resection and chemotherapy.

\subsection{Hirschsprung Disease}
\index{Hirschsprung Disease}
\label{sec:Hirschsprung Disease}
Hirschsprung's disease is a congenital disorder of enteric nervous system development characterized by absence of parasympathetic ganglion cells in the myenteric and submucosal plexuses of the distal colon, with an incidence of 1 in 5,000 live births. The condition results from failure of neural crest cell migration during embryogenesis, most commonly involving the rectosigmoid colon (80\%); long-segment disease extends more proximally. RET proto-oncogene mutations account for 50\% of familial cases. Clinical features include failure to pass meconium within 48 hours of birth (in 90\% of term infants), abdominal distension, bilious vomiting, and chronic constipation. On physical examination, the anal canal is empty and the rectum is tight on digital examination, with explosive expulsion of stool upon withdrawal of the examining finger. Diagnosis is based on rectal biopsy showing absence of ganglion cells and hypertrophied acetylcholinesterase-positive nerve trunks; anorectal manometry shows absence of the rectoanal inhibitory reflex. Management involves surgical resection of the aganglionic segment with pull-through procedure (Swenson, Duhamel, or Soave technique), performed in the neonatal period or after initial colostomy. Prognosis is excellent; most children achieve normal bowel function, though some experience enterocolitis, constipation, or fecal incontinence.

\subsection{Human Metapneumovirus (hMPV)}
\index{Human Metapneumovirus (hMPV)}
\label{sec:pediatric:Human Metapneumovirus}
Human metapneumovirus (hMPV) is a paramyxovirus closely related to respiratory syncytial virus (RSV), discovered in 2001, yet ubiquitous---virtually all children seroconvert by age 5. The condition results from hMPV infecting respiratory epithelium causing syncytia formation, ciliary dysfunction, airway obstruction, and inflammation. Clinical features range from mild URI to severe lower respiratory tract disease (bronchiolitis, croup, pneumonia); severe disease is most common in premature infants, those with cardiopulmonary disease, and immunocompromised children. Diagnosis is by multiplex PCR. Treatment is supportive care; no licensed vaccine or specific antiviral therapy is available. Prognosis is excellent for immunocompetent children.

\subsection{Hypertrophic Pyloric Stenosis}
\index{Hypertrophic Pyloric Stenosis}
\label{sec:Hypertrophic Pyloric Stenosis}
Hypertrophic pyloric stenosis (HPS) is progressive thickening of the pyloric muscle causing gastric outlet obstruction. It occurs in 2--4 per 1000 live births, more commonly in males (4:1) and firstborn infants. The condition results from hypertrophy and hyperplasia of the pyloric circular muscle. Clinical features include presentation at 3--6 weeks of age with non-bilious, projectile vomiting after feeds; infants remain hungry after vomiting. Diagnosis is by physical examination (palpable olive-shaped mass in the right upper quadrant) and ultrasound (thickened pyloric muscle). Metabolic alkalosis with hypochloremia and hypokalemia results from loss of gastric acid. Treatment is pyloromyotomy (Ramstedt procedure) after correction of fluid and electrolyte abnormalities. Prognosis is excellent.

%-------------------------------------------------------------
\section{Genetic and Chromosomal Disorders}
%-------------------------------------------------------------

%-------------------------------------------------------------

%-------------------------------------------------------------

%-------------------------------------------------------------

%-------------------------------------------------------------

%-------------------------------------------------------------

%-------------------------------------------------------------

%-------------------------------------------------------------

\label{sec:Genetic and Chromosomal Disorders}
%-------------------------------------------------------------

\subsection{Intraventricular Hemorrhage (IVH)}
\index{Intraventricular Hemorrhage (IVH)}
\label{sec:Intraventricular Hemorrhage}
Intraventricular hemorrhage (IVH) is bleeding into the germinal matrix and ventricular system of the brain, occurring almost exclusively in preterm infants (<32 weeks gestation). Incidence is 15--25\% in very low birth weight infants. The condition results from fragility of the highly vascularized germinal matrix, which is susceptible to rupture from fluctuations in cerebral blood flow. Clinical features may include apnea, bradycardia, seizures, bulging fontanel, hypotension, and anemia. Diagnosis is by cranial ultrasound. IVH is graded on a scale of I--IV. Management is supportive; prevention is paramount and includes antenatal corticosteroids, delayed cord clamping, and avoidance of rapid volume expansion or hypotension. Prognosis depends on the grade of hemorrhage.

\subsection{Intussusception}
\index{Intussusception}
\label{sec:Intussusception}
Intussusception is the telescoping of one segment of the intestine into an adjacent segment, most commonly ileocecal. It is the most common cause of intestinal obstruction in children aged 3 months to 3 years. The classic triad (not always present) is colicky abdominal pain, vomiting, and red-currant jelly stools (late sign). A sausage-shaped mass may be palpable in the right upper quadrant. Ultrasound demonstrates a target sign or pseudokidney sign. Treatment is non-surgical reduction by air enema under fluoroscopic guidance, which is successful in 80--95\% of cases. Surgical reduction is required for perforation, peritonitis, or failed pneumatic reduction. Prognosis is excellent with early treatment.

\subsection{Kawasaki Disease}
\index{Kawasaki Disease}
\label{sec:pediatric:Kawasaki Disease}
Kawasaki disease is an acute, self-limited vasculitis of unknown etiology that predominantly affects children under 5 years (see \hyperref[ch:connective_tissue]{Chapter~27, Connective Tissue Diseases} for comprehensive coverage). Diagnostic criteria include fever for $\ge$5 days plus four of five features: bilateral conjunctival injection, oral mucosal changes, polymorphous rash, extremity changes, and cervical lymphadenopathy ($\ge$1.5 cm). The most serious complication is coronary artery aneurysms, which develop in 20--25\% of untreated cases. Treatment is intravenous immunoglobulin (IVIG, 2 g/kg) and high-dose aspirin. Prognosis is excellent with prompt treatment; the risk of coronary aneurysms is reduced to less than 5\%.

\subsection{Klinefelter Syndrome (47,XXY)}
\index{Klinefelter Syndrome (47,XXY)}
\label{sec:Klinefelter Syndrome}
Klinefelter syndrome is a chromosomal disorder caused by the presence of an extra X chromosome (47,XXY), with an incidence of 1 in 650 live-born males. Clinical features include tall, eunuchoid habitus, gynecomastia, small testes, azoospermia (infertility), sparse body hair, and low testosterone. Intellectual development is usually normal but verbal and language delays are common. Diagnosis is by karyotype. Management includes testosterone replacement therapy beginning at puberty, breast reduction surgery for gynecomastia, and assisted reproductive technology for infertility.

\subsection{Meconium Aspiration Syndrome}
\index{Meconium Aspiration Syndrome}
\label{sec:Meconium Aspiration Syndrome}
Meconium aspiration syndrome (MAS) occurs when meconium-stained amniotic fluid is aspirated into the fetal airways before or during delivery, causing airway obstruction, pneumonitis, surfactant inactivation, and pulmonary vasoconstriction. It occurs in 1--3\% of all deliveries with meconium-stained fluid. Clinical features include respiratory distress at birth, tachypnea, grunting, retractions, cyanosis, and barrel-shaped chest. Chest radiograph shows coarse, patchy infiltrates, hyperinflation, and atelectasis. Management includes immediate endotracheal suctioning of non-vigorous infants, supplemental oxygen, surfactant therapy, inhaled nitric oxide for pulmonary hypertension, and ECMO for severe refractory cases. The routine practice of intrapartum suctioning for all meconium-stained infants is no longer recommended for vigorous infants. Prognosis is favorable with appropriate management.

\subsection{Necrotizing Enterocolitis (NEC)}
\index{Necrotizing Enterocolitis (NEC)}
\label{sec:Necrotizing Enterocolitis}
Necrotizing enterocolitis (NEC) is a devastating intestinal inflammatory disease of premature infants, with an incidence of 5--10\% in very low birth weight infants. The condition results from a dysregulated inflammatory response to intestinal microbes leading to transmural necrosis, most commonly affecting the terminal ileum and proximal colon. Clinical features include abdominal distension, feeding intolerance, bloody stools, lethargy, and temperature instability. Pneumatosis intestinalis (gas in the bowel wall) on abdominal radiograph is a hallmark sign. Diagnosis is clinical with radiographic confirmation. Treatment includes bowel rest, nasogastric decompression, broad-spectrum antibiotics, and supportive care; surgical intervention is required for perforation or progressive disease. Prognosis is guarded; mortality ranges from 20--50\%.

\subsection{Neonatal Abstinence Syndrome}
\index{Neonatal Abstinence Syndrome}
\label{sec:Neonatal Abstinence Syndrome}
Neonatal abstinence syndrome (NAS) is a postnatal drug withdrawal syndrome in newborns exposed to opioids or other substances in utero. Incidence has risen dramatically with the opioid epidemic. The condition results from chronic fetal exposure to opioids downregulating CNS opioid receptors; abrupt cessation at birth causes neurotransmitter dysregulation and autonomic dysfunction. Clinical features develop within 24--72 hours of birth and include CNS irritability (high-pitched cry, tremors, hypertonia, seizures), autonomic dysfunction (sweating, fever, yawning, sneezing, mottling), and GI dysfunction (poor feeding, vomiting, diarrhea). Diagnosis relies on maternal history with the Finnegan Neonatal Abstinence Scoring System quantifying severity. Treatment: non-pharmacologic care (rooming-in, breastfeeding if mother is on MAT, swaddling, low-stimulation environment) is first-line; pharmacologic therapy includes oral morphine, methadone, or buprenorphine. Prognosis is favorable with treatment; long-term outcomes include increased risk of developmental delays.

\subsection{Neonatal Hypoglycemia}
\index{Neonatal Hypoglycemia}
\label{sec:Neonatal Hypoglycemia}
Neonatal hypoglycemia is defined as plasma glucose <45--50 mg/dL in the first 48 hours of life, though the precise threshold associated with neurodevelopmental impairment remains debated. It occurs in 5--15\% of all newborns and up to 50\% of infants of diabetic mothers. Clinical features are nonspecific: jitteriness, irritability, lethargy, seizures, apnea, poor feeding, and hypothermia. Diagnosis requires timely glucose measurement. Management includes early feeding, intravenous dextrose, and monitoring of glucose levels; persistent or severe hypoglycemia requires evaluation for hyperinsulinism, fatty acid oxidation disorders, and other metabolic diseases. Prognosis is favorable with timely treatment.

\subsection{Neonatal Jaundice / Hyperbilirubinemia / Kernicterus}
\index{Neonatal Jaundice / Hyperbilirubinemia / Kernicterus}
\label{sec:Neonatal Jaundice}
Neonatal jaundice (hyperbilirubinemia) is the yellow discoloration of the skin and sclerae due to elevated serum bilirubin, occurring in 60--80\% of term and nearly all preterm newborns. Physiologic jaundice results from increased bilirubin production, immature hepatic conjugation, and decreased bilirubin excretion. Pathologic jaundice is defined by onset <24 hours of age, rapid rise (>5 mg/dL/day), elevated direct bilirubin, or persistence beyond 2 weeks; causes include hemolytic disease, sepsis, Crigler-Najjar syndrome, Gilbert syndrome, biliary atresia, and breast milk jaundice. Unconjugated bilirubin is neurotoxic; at high levels it crosses the blood-brain barrier and causes bilirubin-induced neurologic dysfunction and kernicterus. Diagnosis includes total and direct serum bilirubin, blood type and Coombs test, G6PD screening, and evaluation for sepsis. Management depends on gestational age, risk factors, and bilirubin level: phototherapy is first-line; exchange transfusion for severe cases or signs of acute bilirubin encephalopathy. Prognosis is favorable with timely treatment.

\subsection{Neonatal Sepsis (Early-Onset and Late-Onset)}
\index{Neonatal Sepsis (Early-Onset and Late-Onset)}
\label{sec:Neonatal Sepsis}
Neonatal sepsis is a systemic bacterial infection in the first 28 days of life. Early-onset sepsis (<72 hours of age) results from vertical transmission from the maternal genital tract, most commonly Group B Streptococcus and Escherichia coli. Late-onset sepsis (>72 hours of age) is nosocomial or community-acquired, with Coagulase-negative staphylococci, S. aureus, E. coli, and Klebsiella being common pathogens. Clinical features are nonspecific: temperature instability, respiratory distress, apnea, bradycardia, feeding intolerance, lethargy, irritability, and seizures. Diagnosis includes blood culture (gold standard), complete blood count, C-reactive protein, and lumbar puncture. Management includes empiric broad-spectrum antibiotics and supportive care; preventative intrapartum antibiotic prophylaxis for GBS-colonized women has significantly reduced early-onset sepsis incidence. Prognosis depends on prompt recognition and treatment.

%-------------------------------------------------------------
\section{Pediatric Cancers}
%-------------------------------------------------------------

%-------------------------------------------------------------

%-------------------------------------------------------------

%-------------------------------------------------------------

Pediatric cancers represent a distinct spectrum of malignancies arising from embryonic cells and mesodermal primitive tissue lines in infants, children, and adolescents. Unlike adult epithelial carcinomas, childhood cancers are dominated by leukemias, brain tumors, lymphomas, and embryonal solid tumors (neuroblastoma, Wilms tumor). Multimodal risk-stratified therapy combines intensive chemotherapy, specialized surgical protocols, and modern radiation strategies to maximize cure rates.

%-------------------------------------------------------------

%-------------------------------------------------------------

%-------------------------------------------------------------

%-------------------------------------------------------------

\label{sec:Pediatric Cancers}
%-------------------------------------------------------------

\subsection{Respiratory Distress Syndrome (RDS)}
\index{Respiratory Distress Syndrome (RDS)}
\label{sec:Respiratory Distress Syndrome}
Respiratory distress syndrome (RDS), also known as hyaline membrane disease, is a pulmonary disorder of prematurity caused by surfactant deficiency in preterm infants. Incidence is inversely proportional to gestational age: approximately 50\% of infants born at 28--30 weeks and nearly all infants born before 26 weeks develop RDS. The condition results from deficiency of surfactant, a phospholipoprotein complex produced by type II pneumocytes that reduces alveolar surface tension. Clinical features include tachypnea, grunting, nasal flaring, intercostal retractions, and cyanosis within minutes to hours after birth. Diagnosis is clinical with chest radiograph showing a characteristic ground-glass appearance with air bronchograms. Prevention includes antenatal corticosteroids. Treatment consists of intratracheal surfactant replacement therapy (beractant, poractant alfa) and respiratory support (nasal continuous positive airway pressure or mechanical ventilation). Prognosis has dramatically improved with surfactant replacement.

\subsection{Retinoblastoma}
\index{Retinoblastoma}
\label{sec:Retinoblastoma}
Retinoblastoma is the most common intraocular malignancy in children, with an incidence of 1 in 15,000--20,000 live births. It arises from primitive retinal cells and is caused by biallelic inactivation of the RB1 tumor suppressor gene. Retinoblastoma occurs in heritable (40\%) and non-heritable (60\%) forms; heritable retinoblastoma involves a germline RB1 mutation with younger age at presentation, bilateral involvement, and increased risk of second malignancies. Clinical features include leukocoria (white pupillary reflex, the most common presenting sign), strabismus, decreased vision, and orbital inflammation. Diagnosis is by ophthalmologic examination under anesthesia. Treatment depends on the extent of disease: intra-arterial chemotherapy, laser photocoagulation, cryotherapy, brachytherapy, systemic chemotherapy, or enucleation. Prognosis is excellent; survival exceeds 95\% in developed countries.

\subsection{Retinopathy of Prematurity (ROP)}
\index{Retinopathy of Prematurity (ROP)}
\label{sec:pediatric:Retinopathy of Prematurity}
Retinopathy of prematurity (ROP) is a vasoproliferative disorder of the developing retina in preterm infants (see \hyperref[ch:eyes]{Chapter~23, Diseases of the Eye} for comprehensive coverage). The condition results from disruption of normal retinal vascularization, with an initial phase of vessel loss followed by pathologic neovascularization. Risk factors include low gestational age (<30 weeks), low birth weight (<1500 g), and supplemental oxygen therapy. Screening by indirect ophthalmoscopy is performed for at-risk preterm infants. Treatment includes laser photocoagulation and intravitreal bevacizumab for severe disease (Type 1 ROP). Prognosis depends on severity and timeliness of treatment.

\subsection{Rhabdomyosarcoma}
\index{Rhabdomyosarcoma}
\label{sec:Rhabdomyosarcoma}
Rhabdomyosarcoma (RMS) is the most common soft tissue sarcoma in children, with an incidence of 4.5 per million children under age 20. It arises from primitive mesenchymal cells with skeletal muscle differentiation. Two major histologic subtypes: embryonal RMS (60\%, more favorable prognosis) and alveolar RMS (20\%, more aggressive). Clinical features depend on the site: a rapidly growing mass, proptosis (orbital), nasal obstruction (parameningeal), or scrotal or vaginal mass. Diagnosis is by biopsy with immunohistochemistry. Treatment is multimodal: chemotherapy, surgical resection, and radiation therapy. Prognosis is 70--80\% survival for low-risk and intermediate-risk, but <50\% for high-risk disease.

\subsection{Rotavirus}
\index{Rotavirus}
\label{sec:Rotavirus}
Rotavirus is the leading cause of severe acute gastroenteritis in infants and young children worldwide before vaccine introduction. Transmission is fecal-oral. The condition results from the virus infecting enterocytes of small intestinal villi, causing cell death, villous atrophy, and malabsorptive diarrhea. Clinical features: incubation period of 1--3 days, acute onset of fever, vomiting followed by profuse watery diarrhea lasting 3--8 days; severe dehydration is the major complication. Diagnosis is clinical; stool antigen detection or PCR is used for confirmation. Treatment: aggressive oral rehydration therapy and zinc supplementation; no antivirals are available. Prevention: live attenuated oral rotavirus vaccines (Rotarix, RotaTeq) have reduced rotavirus hospitalizations by 80--90\%. Prognosis is excellent with appropriate rehydration.

\subsection{Turner Syndrome (45,X)}
\index{Turner Syndrome (45,X)}
\label{sec:Turner Syndrome}
Turner syndrome is a chromosomal disorder resulting from complete or partial monosomy X (45,X, 45,X/46,XX mosaicism, or structural X abnormalities), with an incidence of 1 in 2500 live-born females. Clinical features include short stature, webbed neck, low posterior hairline, shield chest, cubitus valgus, lymphedema of hands and feet at birth, and gonadal dysgenesis (streak gonads, primary amenorrhea, infertility). Associated conditions include coarctation of the aorta, bicuspid aortic valve, horseshoe kidney, sensorineural hearing loss, hypothyroidism, and celiac disease. Intelligence is typically normal. Diagnosis is by karyotype. Management includes recombinant growth hormone therapy for short stature, estrogen replacement for pubertal induction, and ongoing surveillance for associated comorbidities.

\subsection{Viral Exanthems}
\index{Viral Exanthems}
\label{sec:Viral Exanthems}
Viral exanthems are rashes associated with systemic viral infections, most common in childhood. Measles, rubella, and mumps are covered in \hyperref[ch:infectious]{Chapter~20, Infectious Diseases}. Roseola infantum (sixth disease) is caused by human herpesvirus 6 (HHV-6) and less commonly HHV-7, with peak incidence at 6--24 months, presenting with high fever for 3--5 days followed by defervescence and a rose-colored maculopapular rash. Erythema infectiosum (fifth disease) is caused by parvovirus B19, with peak incidence at 5--14 years, presenting with a ``slapped cheek'' facial erythema followed by a lacy reticular rash. Hand-foot-and-mouth disease is caused by enteroviruses, most commonly Coxsackievirus A16 and Enterovirus 71, with peak incidence in children under 5 years, presenting with fever and vesicular rash on hands, feet, and oral mucosa.

\subsection{Williams Syndrome}
\index{Williams Syndrome}
\label{sec:Williams Syndrome}
Williams syndrome (Williams-Beuren syndrome) is a neurodevelopmental disorder caused by a microdeletion on chromosome 7q11.23, encompassing the ELN (elastin) gene. Incidence is 1 in 7500--10,000. Clinical features include distinctive facial appearance (elfin facies), supravalvular aortic stenosis, peripheral pulmonary stenosis, hypercalcemia in infancy, hyperacusis, and a characteristic cognitive profile (strong verbal skills, friendly extroverted personality, but poor visuospatial skills). Diagnosis is by FISH or chromosomal microarray for the 7q11.23 deletion. Management includes cardiac monitoring, dietary calcium restriction during infancy, and educational and behavioral support.

\subsection{Wilms Tumor (Nephroblastoma)}
\index{Wilms Tumor (Nephroblastoma)}
\label{sec:Wilms Tumor}
Wilms tumor is the most common renal malignancy in children, with an incidence of 1 in 10,000 children and median age at diagnosis of 3.5 years. It is associated with genetic syndromes such as WAGR complex, Denys-Drash syndrome, and Beckwith-Wiedemann syndrome. Clinical features include a palpable abdominal mass, abdominal pain, hematuria (15--20\%), and hypertension. Diagnosis is by abdominal ultrasound and CT scan, which shows an intrarenal mass with heterogeneous density. Treatment involves a multidisciplinary approach: neoadjuvant chemotherapy followed by nephrectomy or primary nephrectomy followed by chemotherapy. Prognosis is excellent, with 5-year survival >90\% for localized disease.

%-------------------------------------------------------------
\section{Perinatal Conditions}
%-------------------------------------------------------------

%-------------------------------------------------------------

%-------------------------------------------------------------

%-------------------------------------------------------------

%-------------------------------------------------------------

%-------------------------------------------------------------

%-------------------------------------------------------------

%-------------------------------------------------------------

\label{sec:Perinatal Conditions}
%-------------------------------------------------------------

\subsection{Congenital Diaphragmatic Hernia}
\index{Congenital Diaphragmatic Hernia}
\label{sec:Congenital Diaphragmatic Hernia}
Congenital diaphragmatic hernia (CDH) is the herniation of abdominal contents into the thoracic cavity through a defect in the diaphragm, most commonly the left posterolateral foramen of Bochdalek. Incidence is 1 per 2500 live births. CDH causes pulmonary hypoplasia and pulmonary hypertension, which determine prognosis. Clinical features at birth include respiratory distress, scaphoid abdomen, and decreased breath sounds on the affected side. Diagnosis is confirmed by a nasogastric tube in the chest on radiograph. Management includes immediate endotracheal intubation (avoid bag-mask ventilation), nasogastric decompression, and gentle ventilation strategies; surgical repair is performed after physiologic stabilization. Overall survival is approximately 70--80\%.

\subsection{Congenital Heart Disease (Overview)}
\index{Congenital Heart Disease (Overview)}
\label{sec:Congenital Heart Disease}
Congenital heart disease (CHD) is the most common congenital anomaly, affecting approximately 1\% of live births. CHD encompasses a wide spectrum of structural cardiac defects (see \hyperref[ch:heart]{Chapter~12, Diseases of the Heart} for comprehensive coverage). Septal defects (ASD, VSD, AVSD) allow left-to-right shunting, leading to pulmonary overcirculation and, over time, pulmonary hypertension and Eisenmenger syndrome; small defects may close spontaneously, while larger defects require surgical or transcatheter closure. Tetralogy of Fallot is the most common cyanotic CHD, consisting of VSD, overriding aorta, right ventricular outflow tract obstruction, and right ventricular hypertrophy; surgical repair is performed in infancy. Transposition of the great arteries creates parallel circulations incompatible with life without mixing; the arterial switch operation is performed in the first week of life. Coarctation of the aorta presents with hypertension in the upper extremities and diminished femoral pulses; treatment is surgical resection or balloon angioplasty. Hypoplastic left heart syndrome requires staged surgical palliation (Norwood, Glenn, Fontan procedures) or heart transplantation.

\subsection{Congenital Hip Dysplasia}
\index{Congenital Hip Dysplasia}
\label{sec:Congenital Hip Dysplasia}
Developmental dysplasia of the hip (DDH) encompasses a spectrum from mild acetabular dysplasia to frank dislocation of the femoral head from the acetabulum. Incidence is 1--3 per 1000 live births. Clinical screening includes the Ortolani and Barlow maneuvers. Diagnosis is confirmed by ultrasound (for infants <6 months) or radiography (for older infants). Treatment for infants under 6 months is a Pavlik harness to maintain hip flexion and abduction; for older infants or failed Pavlik treatment, closed or open reduction with a spica cast is required. Prognosis is excellent with early treatment; untreated DDH leads to osteoarthritis and limb-length discrepancy.

\subsection{Congenital Talipes Equinovarus (Clubfoot)}
\index{Congenital Talipes Equinovarus (Clubfoot)}
\label{sec:Congenital Talipes Equinovarus}
Congenital talipes equinovarus (clubfoot) is a congenital deformity of the foot characterized by four components: equinus, varus, adductus, and cavus. Incidence is 1 per 1000 live births, more common in males (2:1), and 50\% are bilateral. The etiology is multifactorial. Diagnosis is primarily clinical at birth. The Ponseti method is the standard of care: serial manipulation and long-leg casting for 6--8 weeks, followed by percutaneous Achilles tenotomy, and then a foot abduction brace. Prognosis is excellent with adherence to bracing; surgical release is reserved for resistant or neglected cases.

\subsection{Cri-du-Chat Syndrome}
\index{Cri-du-Chat Syndrome}
\label{sec:Cri-du-Chat Syndrome}
Cri-du-chat syndrome is a chromosomal disorder caused by deletion of the short arm of chromosome 5 (5p-), with an incidence of 1 in 15,000--50,000 live births. The name derives from the characteristic high-pitched, cat-like cry in infancy due to laryngeal hypoplasia. Clinical features include microcephaly, round face, widely spaced eyes, epicanthal folds, low-set ears, hypotonia, severe intellectual disability, and congenital heart disease. Most cases are de novo deletions. Diagnosis is by karyotype or chromosomal microarray. Management is supportive with early intervention therapies.

\subsection{Croup (Laryngotracheobronchitis)}
\index{Croup (Laryngotracheobronchitis)}
\label{sec:Croup}
Croup is a viral infection of the upper airway causing subglottic edema and obstruction, most commonly caused by parainfluenza virus types 1 and 3. It predominantly affects children 6 months to 3 years of age. Clinical features include abrupt onset of a barking ``seal-like'' cough, inspiratory stridor, hoarseness, and respiratory distress, often worse at night. Diagnosis is clinical; neck radiograph may show the ``steeple sign'' (subglottic narrowing). Treatment includes a single dose of oral or intramuscular dexamethasone (0.6 mg/kg) for mild to moderate croup and nebulized epinephrine for moderate to severe croup. Prognosis is excellent.

\subsection{Cryptorchidism (Undescended Testis)}
\index{Cryptorchidism (Undescended Testis)}
\label{sec:Cryptorchidism}
Cryptorchidism is the absence of one or both testes from the scrotum, occurring in 3\% of full-term and 30\% of preterm male newborns. The condition results from failure of testicular descent. Spontaneous descent occurs in most cases within the first 3--6 months of life. Diagnosis is by physical examination; ultrasound or laparoscopy may be used for nonpalpable testes. Treatment is surgical orchiopexy, optimally performed between 6--18 months of age, to preserve fertility, reduce the risk of testicular torsion, and allow for scrotal examination. Undescended testes are associated with impaired spermatogenesis and increased risk of testicular malignancy; early orchiopexy does not eliminate the cancer risk but facilitates self-examination.

%-------------------------------------------------------------
\section{Prematurity and Its Complications}
%-------------------------------------------------------------

%-------------------------------------------------------------

%-------------------------------------------------------------

%-------------------------------------------------------------

%-------------------------------------------------------------

%-------------------------------------------------------------

%-------------------------------------------------------------

%-------------------------------------------------------------

\label{sec:Prematurity and Its Complications}
%-------------------------------------------------------------

\subsection{Acute Lymphoblastic Leukemia (ALL) in Children}
\index{Acute Lymphoblastic Leukemia (ALL) in Children}
\label{sec:Acute Lymphoblastic Leukemia in Children}
ALL is the most common childhood cancer (25\% of all pediatric malignancies) and the most common leukemia in children (see \hyperref[ch:blood]{Chapter~10, Diseases of the Blood} for comprehensive coverage). Peak incidence is 2--5 years. The majority of childhood ALL is B-cell precursor (B-ALL), and 25\% have the t(12;21) translocation (ETV6-RUNX1 fusion) which confers a favorable prognosis. Clinical features include pallor, fatigue, fever, bone pain, petechiae, bruising, and hepatosplenomegaly. Treatment involves risk-stratified chemotherapy with 5-year survival rates >90\% in children.

\subsection{Acute Otitis Media}
\index{Acute Otitis Media}
\label{sec:Acute Otitis Media in Children}
Acute otitis media (AOM) is the most common bacterial infection in children (see \hyperref[ch:ears]{Chapter~24, Diseases of the Ear} for comprehensive coverage). It affects 80\% of children by age 3, with peak incidence at 6--24 months. Pathogens include Streptococcus pneumoniae, Haemophilus influenzae (nontypeable), and Moraxella catarrhalis. Clinical features include otalgia, fever, irritability, and bulging tympanic membrane with decreased mobility. Diagnosis requires acute onset of symptoms and middle ear effusion on otoscopy. Treatment includes observation (watchful waiting for children >6 months with mild symptoms) or antibiotics (amoxicillin first-line). Prognosis is excellent.

\subsection{Ambiguous Genitalia / Disorders of Sex Development}
\index{Ambiguous Genitalia / Disorders of Sex Development}
\label{sec:Ambiguous Genitalia}
Disorders of sex development (DSD) are congenital conditions in which the development of chromosomal, gonadal, or anatomic sex is atypical. Incidence is approximately 1 in 4500 live births. DSD is classified by karyotype: 46,XX DSD (most commonly congenital adrenal hyperplasia due to 21-hydroxylase deficiency), 46,XY DSD (androgen insensitivity syndrome, 5-alpha-reductase deficiency, gonadal dysgenesis), and sex chromosome DSD (Turner syndrome, Klinefelter syndrome, mixed gonadal dysgenesis). Clinical features include atypical genitalia at birth. Evaluation includes karyotype, hormone levels, imaging, and genitourinary anatomy assessment. Management requires a multidisciplinary team; salt-wasting congenital adrenal hyperplasia is a life-threatening emergency requiring glucocorticoid and mineralocorticoid replacement. Prognosis depends on the specific diagnosis.

\subsection{Angelman Syndrome}
\index{Angelman Syndrome}
\label{sec:Angelman Syndrome}
Angelman syndrome is a neurogenetic disorder caused by loss of function of the maternal copy of the UBE3A gene on chromosome 15q11-q13. Mechanisms include maternal deletion (70\%), paternal uniparental disomy (2--3\%), imprinting defects (3--5\%), and UBE3A mutations (10\%). Incidence is 1 in 12,000--20,000. Clinical features include severe intellectual disability, ataxia, seizures, frequent laughter, happy demeanor, hand-flapping, and absent speech; microcephaly develops after age 2 years. Diagnosis is by methylation analysis of the 15q11-q13 region. Management is supportive, with antiepileptic drugs for seizures and physical/occupational therapy.

\subsection{Brain Tumors in Children}
\index{Brain Tumors in Children}
\label{sec:Brain Tumors in Children}
Brain tumors are the most common solid tumors in children and the second most common pediatric malignancy after leukemia. They differ significantly from adult brain tumors in histology and location. Medulloblastoma is an embryonal tumor of the cerebellum, most common in children aged 3--8 years, presenting with symptoms of increased intracranial pressure, ataxia, and cranial nerve deficits; treatment involves surgical resection, craniospinal radiation, and chemotherapy, with 5-year survival of approximately 70--80\%. Pilocytic astrocytoma is the most common glioma in children, a low-grade (WHO Grade I) tumor typically occurring in the cerebellum, optic pathway, hypothalamus, and brainstem; complete surgical resection is curative, with 90--95\% survival. Ependymoma arises from ependymal cells lining the ventricles, most commonly in the posterior fossa; treatment is maximal safe surgical resection followed by focal radiation, with 5-year survival of 50--70\%. Brainstem glioma, most commonly diffuse intrinsic pontine glioma (DIPG), is a devastating high-grade glioma of the pons in children aged 5--10 years presenting with rapidly progressive cranial neuropathies, ataxia, and long-tract signs; radiation provides temporary palliation, and median survival is less than 1 year.

\subsection{Bronchiolitis (RSV)}
\index{Bronchiolitis (RSV)}
\label{sec:Bronchiolitis}
Bronchiolitis is the most common lower respiratory tract infection in infants, caused most frequently by respiratory syncytial virus (RSV). It affects nearly all children by age 2, with peak incidence at 2--6 months. The condition results from inflammation, edema, and necrosis of the bronchiolar epithelium, leading to airway obstruction, air trapping, and atelectasis. Clinical features include coryza, cough, tachypnea, wheezing, crackles, and respiratory distress; apnea is a presenting sign in young infants (<2 months). Diagnosis is clinical; RSV rapid antigen testing or PCR confirms the etiology. Treatment is primarily supportive: supplemental oxygen, nasal suctioning, and hydration. Palivizumab (RSV monoclonal antibody) is used for prophylaxis in high-risk infants. Prognosis is excellent with supportive care.

\subsection{Bronchopulmonary Dysplasia (BPD)}
\index{Bronchopulmonary Dysplasia (BPD)}
\label{sec:Bronchopulmonary Dysplasia}
Bronchopulmonary dysplasia (BPD) is a chronic lung disease of prematurity, defined as the need for supplemental oxygen or respiratory support at 36 weeks postmenstrual age. Incidence is 30--50\% in extremely preterm infants. The condition results from the combined effects of pulmonary immaturity, mechanical ventilation, oxygen toxicity, and inflammation, leading to impaired alveolarization and dysregulated vascular development. Clinical features include persistent oxygen requirement, tachypnea, retractions, and wheezing. Diagnosis is clinical. Management includes optimizing nutrition (high-calorie feeds), fluid restriction, diuretics, bronchodilators, and systemic corticosteroids in severe cases. Prognosis is variable; long-term complications include reactive airway disease, pulmonary hypertension, and increased risk of respiratory infections.

\subsection{Cleft Lip and Palate}
\index{Cleft Lip and Palate}
\label{sec:Cleft Lip and Palate}
Cleft lip and palate are among the most common congenital anomalies (1 per 700 live births), resulting from failure of fusion of the maxillary and medial nasal processes (cleft lip) or the palatal shelves (cleft palate). Clinical features include feeding difficulties, speech abnormalities, hearing loss, and dental anomalies. Diagnosis is clinical at birth. Treatment requires a multidisciplinary team approach: surgical repair (cheiloplasty at 3--6 months, palatoplasty at 9--18 months), speech therapy, orthodontic intervention, and management of associated otologic issues. Prognosis is excellent with appropriate multidisciplinary care.

